\chapter{Quantum Basics}

\section{Second Quantization} 

\subsection{One-particle basis functions}

Any basis of a space must span the space and consist of linearly
independent elements of that space. The formulation here is restricted
to finite-dimensional vector/function spaces equipped with an inner
product. In practice, this will be the case anyway, and the
formulation is simpler this way.
We can formulate the condition of $\{|\varphi_n\rangle\}$ spanning the
space $\Phi$ by requiring, for all $|\varphi\rangle\in\Phi$, that there
exists a set of coefficients $\{c_1,\dots,c_N\}$ such that

\begin{equation}\label{eq:basis-expansion}
|\varphi\rangle = \sum_{n=1}^N c_n|\varphi_n\rangle = \big(|\varphi_1\rangle\ \cdots\ |\varphi_N\rangle\big)\begin{pmatrix}c_1\\\vdots\\c_N\end{pmatrix}.
\end{equation}

We can find a linear system of equations for the coefficients by
considering

\begin{equation}
\langle\varphi_m|\varphi\rangle = \sum_{n=1}^N c_n \langle\varphi_m|\varphi_n\rangle,
\end{equation}

so the linear system of equations becomes

\begin{equation}\label{eq:gram-matrix}
\begin{pmatrix}\langle\varphi_1|\varphi\rangle\\\vdots\\\langle\varphi_N|\varphi\rangle\end{pmatrix}
=
\begin{pmatrix}\langle\varphi_1|\varphi_1\rangle & \cdots & \langle\varphi_1|\varphi_N\rangle\\ \vdots & & \vdots \\ \langle\varphi_N|\varphi_1\rangle & \cdots & \langle\varphi_N|\varphi_N\rangle\end{pmatrix}
\begin{pmatrix}c_1\\\vdots\\c_N\end{pmatrix},
\end{equation}

with the Gram matrix $\mathbf{S}$:

\begin{equation}
\mathbf{S} = \begin{pmatrix}\langle\varphi_1|\varphi_1\rangle & \cdots & \langle\varphi_1|\varphi_N\rangle\\ \vdots & & \vdots \\ \langle\varphi_N|\varphi_1\rangle & \cdots & \langle\varphi_N|\varphi_N\rangle\end{pmatrix}.
\end{equation}

Solving for $\vec c$, we find

\begin{equation}
\vec c = \mathbf{S}^{-1}\begin{pmatrix}\langle\varphi_1|\varphi\rangle\\\vdots\\\langle\varphi_N|\varphi\rangle\end{pmatrix},
\end{equation}

and inserting $\vec c$ back into \eqref{eq:basis-expansion}, we find

\begin{equation}\label{eq:completeness-eigenfunction-phi}
|\varphi\rangle = \sum_n\sum_m |\varphi_n\rangle\left(\mathbf{S}^{-1}\right)_{nm}\langle\varphi_m|\varphi\rangle
\end{equation}

and therefore

\begin{equation}\label{eq:completeness-eigenfunction}
\sum_n\sum_m |\varphi_n\rangle\left(\mathbf{S}^{-1}\right)_{nm}\langle\varphi_m| = \mathbb{1}.
\end{equation}

\subsubsection*{Proving the Gram matrix is invertible}

$\mathbf{S}$ invertible implies, for a finite-dimensional space, that
the kernel must be trivial:

\begin{equation}
\mathbf{S}\vec a = 0 \iff \vec a = 0.
\end{equation}

Consider

\begin{equation}
\begin{pmatrix}\langle\varphi_1|\varphi_1\rangle & \cdots & \langle\varphi_1|\varphi_N\rangle\\ \vdots & & \vdots \\ \langle\varphi_N|\varphi_1\rangle & \cdots & \langle\varphi_N|\varphi_N\rangle\end{pmatrix}
\begin{pmatrix}a_1\\\vdots\\a_N\end{pmatrix}
=
\begin{pmatrix}0\\\vdots\\0\end{pmatrix}.
\end{equation}

Then, for all $m=1,\dots,N$,

\begin{equation}
\Big\langle\varphi_m\Big|\sum_{n=1}^N a_n\varphi_n\Big\rangle = 0
\end{equation}

holds. In words: there is a vector $|\varphi_a\rangle=\sum_{n=1}^N a_n|\varphi_n\rangle$
which lies in $\Phi$ and is yet orthogonal to each vector
$|\varphi_m\rangle$, $m=1,\dots,N$, which span the space and thereby also to itself. The only
vector that fulfils this requirement is the zero vector. Therefore we
find the new condition:

\begin{equation}
\sum_{n=1}^N a_n |\varphi_n\rangle \overset{!}{=} 0.
\end{equation}

Since the vectors of the set $\{|\varphi_n\rangle\}_{n=1}^N$ are linearly independent
, this equation is only satisfied by $\vec a=\vec 0$. Meaning
that the kernel is trivial, which implies $\mathbf{S}$ is invertible.

\subsubsection*{Going to infinite dimensions} 
In the infinite-dimensional case the requirements have
to be formulated more carefully. The kernel being trivial will not
suffice for the operator $\mathbf{S}$ to be invertible: since the
operator is now a map between two countably infinite-dimensional
spaces, the infinite-matrix picture for $\mathbf{S}$ remains valid, but
$\mathbf{S}$ must additionally fulfil the property that its spectrum be
bounded from above and below. In the finite-dimensional case,
boundedness from below and trivial kernel are equivalent. An
infinite-dimensional example of a non-invertible operator with trivial
kernel is

\begin{equation}
\mathbf{S} = \begin{pmatrix} 1 & & & 0 \\ & \tfrac12 & & \\ & & \tfrac13 & \\ 0 & & & \ddots \end{pmatrix}.
\end{equation}

The kernel is trivial, but there exists no $A>0$ such that for all
$c\in\ell^2$,

\begin{equation}
A\|c\| \le \|\mathbf{S}c\|.
\end{equation}

This means the range that $\mathbf{S}$ maps onto is \emph{not closed},
and the inverse operator has diverging entries.

To see the problem with unbounded operators like $\mathbf{S}^{-1}$ in this
example:

\begin{equation}
\mathbf{S}^{-1} = \begin{pmatrix} 1 & & & \\ & 2 & & \\ & & 3 & \\ & & & \ddots \end{pmatrix}
\end{equation}

Consider the sequence $a_n = (1/n)_{n\in\mathbb{N}}$, which lies within
$\ell^2$ since $\sum_{n\in\mathbb{N}}(1/n)^2$ converges. However,

\begin{equation}
\mathbf{S}^{-1}a = (1,1,1,\dots),
\end{equation}

where $\sum_{n\in\mathbb{N}}1^2\to\infty$ is not finite. The operator
$\mathbf{S}^{-1}$ therefore maps a sequence that lies within $\ell^2$
out of $\ell^2$.

\subsubsection*{The bra is not the dual element}

Let $\{|\varphi_n\rangle\}$ be the set of basis vectors of a vector
space $V$. Then the elements of the dual set $\{\langle\tilde\varphi_n|\}$
have the property:

\begin{equation}
\langle\tilde\varphi_m|\varphi_n\rangle = \delta_{mn}.
\end{equation}

This property is obviously not fulfilled by the definition of the
bra's (only for orthogonal vectors). We can find the dual elements via
equation \eqref{eq:identity-eigenfunction} to be:

\begin{equation}
\langle\tilde\varphi_m| = \sum_{m'}\left(\mathbf{S}^{-1}\right)_{mm'}\langle\varphi_{m'}|.
\end{equation}

The bra $\langle\varphi_m|$ is still an element of the dual space,
however with:

\begin{align}
\langle\varphi_m|&: V\to\mathbb K\\
|\varphi\rangle &\mapsto \int dx\,\varphi_m^*(x)\,\varphi(x).\label{eq:functional-map-to-inner}
\end{align}

For an orthogonal basis, this definition agrees with the definition of
the dual element, for that reason it often goes unnoticed that $\langle\varphi_m|$ is
not in general the dual element to $|\varphi_m\rangle$.

\subsection*{The position and momentum bra's and ket's}

Consider the definition of the scalar product of two wave functions as
demonstrated in \eqref{eq:functional-map-to-inner}. We can interpret $\varphi(x)$
as the result of a linear functional $X_{LF}$ mapping $\varphi$ to its
value $\varphi(x)$ at $x$:

\begin{align}
X_{LF}:\ \mathcal S &\to \mathbb C\\
\varphi &\mapsto \varphi(x)
\end{align}

$\mathcal S$ just denotes the physical subspace of the Hilbert space
$L^2$ of square-integrable functions \cite{Madrid_2005}. Similarly, we can interpret
$\varphi_m^*(x)$ as the result of the anti-linear functional $X_{AF}$
mapping $\varphi_m$ to the complex conjugate of its value
$\varphi_m^*(x)$ at $x$:

\begin{align}
X_{AF}:\ \mathcal S &\to \mathbb C\\
\varphi &\mapsto \varphi^*(x)
\end{align}

With this interpretation we can rewrite the inner product:

\begin{equation}
\langle\varphi|\psi\rangle = \int dx\,\varphi^*(x)\psi(x) = \int dx\,X_{AF}(\varphi)\,X_{LF}(\psi). \label{eq:inner-product-pos-rep}
\end{equation}

Now we define the ket $|x\rangle$ to be the anti-linear functional and
the bra $\langle x|$ to be the linear functional, since the inner
product here is defined to be anti-linear in the first entry and
linear in the second entry. Therefore we can rewrite
\eqref{eq:inner-product-pos-rep} to:

\begin{equation}
\langle\varphi|\psi\rangle = \int dx\,\langle\varphi|x\rangle\langle x|\psi\rangle.
\end{equation}

Further we find, if $\varphi(x)$ and $\psi(x)$ are connected via a
Fourier transform with $\varphi(p)$ and $\psi(p)$ respectively, then

\begin{equation}
\int dx\,\varphi^*(x)\psi(x) = \int dp\,\varphi^*(p)\psi(p).
\end{equation}

\begin{proof}
\begin{align}
\int dx\,\varphi^*(x)\psi(x)
&= \int dx\int dp\,\frac{1}{\sqrt{2\pi\hbar}}e^{-ipx/\hbar}\varphi^*(p)
\int dp'\,e^{ip'x/\hbar}\frac{1}{\sqrt{2\pi\hbar}}\psi(p') \\
&= \int dp\,dp'\,\varphi^*(p)\psi(p')\,\frac{1}{2\pi\hbar}\int dx\,e^{ix(p'-p)/\hbar}\\
&= \int dp\,dp'\,\varphi^*(p)\psi(p')\,\delta(p-p')\\
&= \int dp\,\varphi^*(p)\psi(p).
\end{align}
\end{proof}

So with the same arguments used for the position bra and ket, we find
for momentum:

\begin{equation}
\langle\varphi|\psi\rangle = \int dp\,\langle\varphi|p\rangle\langle p|\psi\rangle.
\end{equation}

If we consider the Fourier relationship between $\varphi(x)$ and
$\varphi(p)$, we can find an additional interpretation of the position
and momentum bra's as plane waves:

\begin{equation}
\langle x|\varphi\rangle = \varphi(x) = \int dp\,\frac{1}{\sqrt{2\pi\hbar}}e^{ipx/\hbar}\varphi(p).
\end{equation}

At the same time we can interpret the scalar product over two
momentum-dependent functions directly:

\begin{equation}
\langle x|\varphi\rangle = \int dp\,x^*(p)\varphi(p) = \int dp\,\langle x|p\rangle\langle p|\varphi\rangle,
\end{equation}

so we find

\begin{equation}
\langle p|x\rangle = x(p) = \frac{1}{\sqrt{2\pi\hbar}}e^{-ipx/\hbar}
\end{equation}

and analogously

\begin{equation}
p(x) = \langle x|p\rangle = \langle p|x\rangle^* = \frac{1}{\sqrt{2\pi\hbar}}e^{ipx/\hbar}.
\end{equation}

We can also use the plane-wave interpretation to determine the
overlaps $\langle x|x'\rangle$ and $\langle p|p'\rangle$ by directly
computing the inner product of the plane-wave functions:

\begin{align}
\langle x|x'\rangle &= \int dp\,\frac{1}{2\pi\hbar}e^{ipx/\hbar}e^{-ipx'/\hbar} = \delta(x-x')\\
\langle p|p'\rangle &= \int dx\,\frac{1}{2\pi\hbar}e^{-ipx/\hbar}e^{ip'x/\hbar} = \delta(p-p').
\end{align}

At this point I should state that there are also approaches \cite{Madrid_2005}
which determine various "inner products" of functionals purely as
solutions to differential equations. Also a very nice approach, which remains loyal to the functional interpretation.
The only downside is, that one is stuck with many bra-kets which are not understood as inner products.
\paragraph{Note:}
That being said the careful reader probably realized that the bra-ket $\langle x|p\rangle$ and its complex conjugate
can also not be understood as an inner product. In this derivation we find its value via the fourier relations. Also 
the interpretation of the position and momentum bra's and ket's as functionals is a preliminary to finding the inner product 
interpretation later. One can also interpret the plane-waves as kernels to the functionals.

\subsection*{The completeness relation in position and momentum representation}

In the following, $\alpha$ can represent position $x$ or momentum $p$.
We consider the completeness relation \eqref{eq:completeness-eigenfunction-phi} in
$\alpha$-representation and expand the inner product in the
$\alpha$-representation:

\begin{align}
\langle\alpha|\varphi\rangle &= \sum_n\sum_m \langle\alpha|\varphi_n\rangle\left(\mathbf S^{-1}\right)_{nm}\int d\alpha'\,\langle\varphi_m|\alpha'\rangle\langle\alpha'|\varphi\rangle\\
\varphi(\alpha) &= \sum_n\sum_m \varphi_n(\alpha)\left(\mathbf S^{-1}\right)_{nm}\int d\alpha'\,\varphi_m^*(\alpha')\varphi(\alpha')
\end{align}

So we find, within an integral over $\alpha'$:

\begin{equation}
\sum_n\sum_m \varphi_n(\alpha)\left(\mathbf S^{-1}\right)_{nm}\varphi_m^*(\alpha') = \delta(\alpha-\alpha').
\end{equation}

We can also consider the inner product in momentum representation

\begin{align}
\langle x|\varphi\rangle &= \sum_n\sum_m \langle x|\varphi_n\rangle\left(\mathbf S^{-1}\right)_{nm}\int dp\,\langle\varphi_m|p\rangle\langle p|\varphi\rangle\\
&= \int dp \sum_n\sum_m \varphi_n(x)\left(\mathbf S^{-1}\right)_{nm}\varphi_m^*(p)\,\varphi(p)
\end{align}

and find the mixed representation completeness relation

\begin{equation}
\sum_n\sum_m \varphi_n(x)\left(\mathbf S^{-1}\right)_{nm}\varphi_m^*(p) = \frac{1}{\sqrt{2\pi\hbar}}\exp(ipx/\hbar).
\end{equation}

\subsection{Orbital Creation and Annihilation operators}
It is assumed that the diagonal entries of the Gram-Matrix are normalized. Ensuring normalization is practically not a problem
and without normalized states the formulas become unnecessarily crowded.

\paragraph{Bosons:}
\begin{equation}
  [\hat b_i, \hat b_j] = 0, \qquad
  [\hat b^{\dagger}_i, \hat b^{\dagger}_j] = 0
\end{equation}
\textbf{Case 1:}\quad $[\hat b_i, \hat b_j^\dagger] = \delta_{ij}$
\begin{equation}
|n_1,\dots,n_L\rangle = \prod_{k=1}^L \frac{\left(\hat b_k^\dagger\right)^{n_k}}{\sqrt{n_k!}}\,|\mathrm{vac}\rangle
\end{equation}

\textbf{Case 2:}\quad $[\hat b_i, \hat b_j^\dagger] = \mathbf S_{ij}$
\begin{equation}
|n_1,\dots,n_L\rangle = \frac{1}{\sqrt{\operatorname{perm}(\mathbf S_{\mathrm{occ}})}}\prod_{k=1}^L \left(\hat b_k^\dagger\right)^{n_k}\,|\mathrm{vac}\rangle
\end{equation}

The $\dfrac{1}{\sqrt{n_k!}}$ term is just a special case of the more general 
permanent formula (The permanent is contains the same terms as the determinant but all signs are plus signs).  
The factor ensures normalization of the state. This can
easily be shown just using the commutation relations of the operators.

\paragraph{Fermions:}
\begin{equation}
  \{\hat c^{\dagger}_i, \hat c^{\dagger}_j\} = 0, \qquad
  \{\hat c_i, \hat c_j\} = 0, \qquad
  \{\hat c_i, \hat c^{\dagger}_j\} = \mathbf{S}_{ij}.
\end{equation}

\paragraph{First-quantized/Orbital-list:}
\begin{equation}
  |\vec k\rangle = |k_1, \dots, k_N\rangle
  = \frac{1}{\sqrt{\operatorname{det}(\mathbf S_{\mathrm{occ}})}}\prod_{i=1}^{N} \hat c^{\dagger}_{k_i}\,|\text{vac}\rangle.
\end{equation}
More intuitive and useful for theoretical considerations.

\paragraph{Second-quantized/Occupation-number:}
\begin{align}
  \hat c^{\dagger}_i\,|n_1, \dots, n_i, \dots, n_L\rangle
  &= (-1)^{\sum_{j=1}^{i-1} n_j}\,(1 - n_i)\,
     |n_1, \dots, n_i + 1, \dots, n_L\rangle, \\[4pt]
  \hat c_i\,|n_1, \dots, n_i, \dots, n_L\rangle
  &= (-1)^{\sum_{j=1}^{i-1} n_j}\,n_i\,
     |n_1, \dots, n_i - 1, \dots, n_L\rangle.
\end{align}
This notation is used in computational approaches. Here the normalization factor of $\frac{\det(\mathbf S_{\mathrm{occ,\,after}})}{\det(\mathbf S_{\mathrm{occ,\,before}})}$ is 
neglected.

\subsection*{The mixed anti-commutator for fermions}

Postulating the existence of the anti-commutator immediately returns its value:

\begin{align*}
\langle\varphi_n|\varphi_{n'}\rangle &= \langle\mathrm{vac}|\hat c_n \hat c_{n'}^\dagger|\mathrm{vac}\rangle \\
&= -\langle\mathrm{vac}|\hat c_{n'}^\dagger \hat c_n|\mathrm{vac}\rangle + \{\hat c_n, \hat c_{n'}^\dagger\}
\end{align*}

So we find

\begin{equation}
\{\hat c_n, \hat c_{n'}^\dagger\} = \langle\varphi_n|\varphi_{n'}\rangle = \left(\mathbf{S}\right)_{nn'}
\end{equation}

We can also see that this definition returns the expected overlaps:

\begin{align*}
\langle\varphi_n\varphi_m|\varphi_{n'}\varphi_{m'}\rangle &= \langle\mathrm{vac}|\hat c_n\hat c_m \hat c_{m'}^\dagger\hat c_{n'}^\dagger|\mathrm{vac}\rangle \\
&= -\langle\mathrm{vac}|\hat c_n \hat c_{m'}^\dagger \hat c_m \hat c_{n'}^\dagger|\mathrm{vac}\rangle + \langle\mathrm{vac}|\hat c_n\hat c_{n'}^\dagger|\mathrm{vac}\rangle\,\langle\varphi_m|\varphi_{m'}\rangle \\
&= -\langle\mathrm{vac}|\hat c_n\hat c_{m'}^\dagger|\mathrm{vac}\rangle\,\langle\varphi_m|\varphi_{n'}\rangle + \langle\varphi_n|\varphi_{n'}\rangle\langle\varphi_m|\varphi_{m'}\rangle \\
&= -\langle\varphi_n|\varphi_{m'}\rangle\langle\varphi_m|\varphi_{n'}\rangle + \langle\varphi_n|\varphi_{n'}\rangle\langle\varphi_m|\varphi_{m'}\rangle
\end{align*}

For the Slater determinants we would find

\begin{align*}
\int dx\,dy\ &\frac{1}{\sqrt2}\Big(\varphi_n(x)\varphi_m(y) - \varphi_m(x)\varphi_n(y)\Big)^*
\frac{1}{\sqrt2}\Big(\varphi_{n'}(x)\varphi_{m'}(y) - \varphi_{m'}(x)\varphi_{n'}(y)\Big) \\
&= \frac{1}{2}\Big[\langle\varphi_n|\varphi_{n'}\rangle\langle\varphi_m|\varphi_{m'}\rangle - \langle\varphi_n|\varphi_{m'}\rangle\langle\varphi_m|\varphi_{n'}\rangle
- \langle\varphi_n|\varphi_{m'}\rangle\langle\varphi_m|\varphi_{n'}\rangle + \langle\varphi_n|\varphi_{n'}\rangle\langle\varphi_m|\varphi_{m'}\rangle \Big] \\
&= \langle\varphi_n|\varphi_{n'}\rangle\langle\varphi_m|\varphi_{m'}\rangle - \langle\varphi_n|\varphi_{m'}\rangle\langle\varphi_m|\varphi_{n'}\rangle
\end{align*}

If the sets $\{n,m\}$ and $\{n',m'\}$ contain the same
elements/orbitals, say $n=n'$ and $m=m'$, we find:

\begin{align}
\langle\varphi_n\varphi_m|\varphi_n\varphi_m\rangle
&= \langle\varphi_n|\varphi_n\rangle\langle\varphi_m|\varphi_m\rangle - \langle\varphi_n|\varphi_m\rangle\langle\varphi_m|\varphi_n\rangle\\
&= \det(\mathbf S_{\mathrm{occ}})
\end{align}

where the subscript ``occ'' denotes that $\mathbf S_{\mathrm{occ}}$ is
the ``small'' Gram matrix of the occupied orbitals, not the one that
contains the overlaps between all basis states:

\begin{equation}
\mathbf S_{\mathrm{occ}} = \begin{pmatrix}\langle\varphi_n|\varphi_n\rangle & \langle\varphi_n|\varphi_m\rangle\\ \langle\varphi_m|\varphi_n\rangle & \langle\varphi_m|\varphi_m\rangle\end{pmatrix}.
\end{equation}

If we had chosen $n=m'$ and $m=n'$, we would have found an additional
minus sign, corresponding to $\mathrm{sgn}(\sigma)=-1$ for the
transposition.

To visualize the prefactor of each term, we first start talking about
orbitals, because the annihilation operators correspond to bra-orbitals
and the creation operators to ket-orbitals. Standard-ordering is
achieved if we take the overlap of adjacent bra and ket orbitals always:

\begin{align*}
\langle\varphi_n\varphi_m|\varphi_{n'}\varphi_{m'}\rangle
&= \langle\mathrm{vac}|\hat c_n\hat c_m\hat c_{m'}^\dagger\hat c_{n'}^\dagger|\mathrm{vac}\rangle \notag\\
&\big\Downarrow \notag\\
&\Big[(\langle\varphi_n|\langle\varphi_m|)(|\varphi_{m'}\rangle|\varphi_{n'}\rangle)\Big] \notag\\
&\quad \langle\varphi_m|\varphi_{m'}\rangle\ \big[\langle\varphi_n|\big]\big[|\varphi_{n'}\rangle\big] \notag\\
&\quad \langle\varphi_n|\varphi_{n'}\rangle\langle\varphi_m|\varphi_{m'}\rangle
\end{align*}

Whenever one anti-commutes with an operator/orbital first, before
evaluating the anticommutator, that corresponds to an orbital appearing
to the \emph{left} of the orbital even though it should be to the
right in standard-ordering. So the prefactor becomes $(-1)^{n_L}$,
where $n_L$ denotes the orbitals that are matched with an "earlier" (further to the left)
bra-orbital than they would be in normal ordering --- or equivalently,
the number of ket-orbitals that should be to the right of the
corresponding ket-orbital. The product of all those factors represents
the permutation-factor, as the number of permutations needed
corresponds to the sum of the $n_L$.

\textbf{Proof.} That the number of permutations corresponds to the sum
of $n_L$ for each position.

\emph{Strategy:} Start with the left-most position with $n_L > 0$ and
commute the orbital through its left neighbors. The orbitals that
$n_L$ refers to must be direct neighbors, otherwise you didn't choose
the left-most position/orbital. Now iterate the procedure until you
have found standard-ordering.

Now we can also quickly convince ourselves that we find the same
pattern for the overlap of two Slater determinants. First, for the
bra Slater determinant we only consider the normal-ordered term:

\begin{equation*}
\varphi_1(x_1)\cdots\varphi_N(x_N)
\end{equation*}

Therefore the sign can only come from the term of the Slater
determinant and corresponds to $(-1)^n$, with $n$ being the number of
permutations required to reach normal-ordering. For the other bra
terms we use the trick:

\begin{align*}
\langle\varphi_1|\varphi_1'\rangle \cdots \langle\varphi_j|\varphi_i'\rangle \cdots \langle\varphi_i|\varphi_j'\rangle \cdots \langle\varphi_N|\varphi_N'\rangle
&= \langle\varphi_1|\varphi_1'\rangle \cdots \langle\varphi_i|\varphi_j'\rangle \cdots \langle\varphi_j|\varphi_i'\rangle \cdots \langle\varphi_N|\varphi_N'\rangle
\end{align*}

where $i<j$. In other words, we can move the permutations of the bra
term onto the ket term, and we again find that the sign is just
$(-1)^n$. So the behaviour is the same.

\textbf{Note:} Here, in the Slater-determinant case, one finds all possible combinations  
for each of the $N!$ bra terms; Also, notice that for the overlap of a Slater determinant 
with itself all possible combinations just become $\det(\mathbf{S_{occ}})$ where $\mathbf{S_{occ}}$ is built with all the 
occupied orbitals. however, through
the normalization of bra- and ket-state we find the factor $\tfrac{1}{N!\det(\mathbf{S_{occ}})}$.

\subsection*{Field Operators}

\paragraph{Fundamental property:}
\begin{equation}
\hat\varPsi^\dagger(\alpha)|\mathrm{vac}\rangle = |\alpha\rangle
\qquad\text{and}\qquad
\langle\mathrm{vac}|\hat\varPsi(\alpha) = \langle\alpha|,
\end{equation}
where $\alpha\in\{x,p\}$.

We can use \eqref{eq:completeness-eigenfunction}, for which we demonstrated that it
holds between any bra-ket $\langle\gamma|\dots|\gamma'\rangle$
with $\gamma,\gamma'\in\{x,p,\varphi_k\}$ (For the specific case of the eigenfunctions $\{\phi_{k}\}$ that \eqref{eq:completeness-eigenfunction}
is an identity just acting on an element of the function space spanned by the eigenfunctions \eqref{eq:completeness-eigenfunction-phi}). 
Therefore we can expand the
anti-linear functional $|\alpha\rangle$ and evaluate against the
$\{\varphi_k\}$, which span the space of test functions:

\begin{align}
\langle\varphi_k|\hat\varPsi^\dagger(\alpha)|\mathrm{vac}\rangle
&= \langle\varphi_k|\sum_n\sum_m|\varphi_n\rangle\left(\mathbf S^{-1}\right)_{nm}\langle\varphi_m|\alpha\rangle\\
\langle\varphi_k|\hat\varPsi^\dagger(\alpha)|\mathrm{vac}\rangle
&= \langle\varphi_k|\sum_n\sum_m \hat c_n^\dagger|\mathrm{vac}\rangle\left(\mathbf S^{-1}\right)_{nm}\varphi_m^*(\alpha)
\end{align}

Since both functionals agree on all basis vectors, they agree on the
entire space of test functions. And since the operators on both sides
consist only of creation operators, we find operator equivalence:

\begin{equation}
\boxed{\hat\varPsi^\dagger(\alpha) = \sum_n\sum_m\left(\mathbf S^{-1}\right)_{nm}\varphi_m^*(\alpha)\,\hat c_n^\dagger}
\end{equation}

Before taking the hermitian conjugate, note that $\mathbf S^{-1}$ is
hermitian, since $\mathbf S$ is hermitian:

\begin{equation}
\mathbf S\,\mathbf S^{-1} = \mathbb 1
\ \Longrightarrow\ \left(\mathbf S^{-1}\right)^\dagger\mathbf S^\dagger = \left(\mathbf S^{-1}\right)^\dagger\mathbf S = \mathbb 1
\ \Longrightarrow\ \left(\mathbf S^{-1}\right)^\dagger = \mathbf S^{-1},
\end{equation}

so that $\left(\mathbf S^{-1}\right)^*_{nm} = \left(\mathbf S^{-1}\right)_{mn}$.
Taking the hermitian conjugate of the operator equivalence and
relabelling $n\leftrightarrow m$:

\begin{equation}
\boxed{\hat\varPsi(\alpha) = \sum_n\sum_m\left(\mathbf S^{-1}\right)_{nm}\varphi_n(\alpha)\,\hat c_m}
\end{equation}

Further, we found that $\int d\alpha\,|\alpha\rangle\langle\alpha|$ acts
like an identity inside any bra-ket $\langle\gamma|\dots|\gamma'\rangle$
with $\gamma,\gamma'\in\{x,p,\varphi_k\}$. We use this to expand the
eigenfunctions and again test against the eigenfunctions:

\begin{align}
\langle\varphi_k|\hat c_n^\dagger|\mathrm{vac}\rangle &= \langle\varphi_k|\int d\alpha\,|\alpha\rangle\langle\alpha|\hat c_n^\dagger|\mathrm{vac}\rangle\\
\langle\varphi_k|\hat c_n^\dagger|\mathrm{vac}\rangle &= \langle\varphi_k|\int d\alpha\,\hat\varPsi^\dagger(\alpha)|\mathrm{vac}\rangle\,\varphi_n(\alpha)
\end{align}

With the same arguments as before we find the operator identity:

\begin{equation}
\boxed{\hat c_n^\dagger = \int d\alpha\,\hat\varPsi^\dagger(\alpha)\,\varphi_n(\alpha)}
\end{equation}

and by taking the hermitian conjugate

\begin{equation}
\boxed{\hat c_n = \int d\alpha\,\varphi_n^*(\alpha)\,\hat\varPsi(\alpha)}
\end{equation}

\paragraph{Consistency check.} Substituting the field expansion into
the mode expansion returns the mode operator:

\begin{equation}
\int d\alpha\,\hat\varPsi^\dagger(\alpha)\,\varphi_n(\alpha)
= \sum_{k,m}\left(\mathbf S^{-1}\right)_{km}\hat c_k^\dagger
\underbrace{\int d\alpha\,\varphi_m^*(\alpha)\varphi_n(\alpha)}_{=\,S_{mn}}
= \sum_k \delta_{kn}\,\hat c_k^\dagger = \hat c_n^\dagger.
\end{equation}

With the relationship between the field operators and
creation/annihilation operators, we can establish the
anti-commutators:

\begin{align}
\{\hat\varPsi(\alpha),\hat\varPsi^\dagger(\alpha')\}
&= \sum_n\sum_m\sum_{n'}\sum_{m'}\left(\mathbf S^{-1}\right)_{nm}\varphi_n(\alpha)\left(\mathbf S^{-1}\right)_{n'm'}\varphi_{m'}^*(\alpha')\,\{\hat c_m,\hat c_{n'}^\dagger\}\\
&= \sum_n\sum_m\sum_{n'}\sum_{m'}\left(\mathbf S^{-1}\right)_{nm}\left(\mathbf S\right)_{mn'}\left(\mathbf S^{-1}\right)_{n'm'}\varphi_n(\alpha)\varphi_{m'}^*(\alpha')\\
&= \sum_n\sum_{m'}\left(\mathbf S^{-1}\right)_{nm'}\varphi_n(\alpha)\varphi_{m'}^*(\alpha')\\
&= \delta(\alpha-\alpha')
\end{align}

and for the mixed anti-commutators:

\begin{align}
\{\hat\varPsi(\alpha),\hat c_k^\dagger\}
&= \sum_n\sum_m\left(\mathbf S^{-1}\right)_{nm}\varphi_n(\alpha)\,\{\hat c_m,\hat c_k^\dagger\}\\
&= \sum_n\sum_m\left(\mathbf S^{-1}\right)_{nm}\left(\mathbf S\right)_{mk}\varphi_n(\alpha)\\
&= \sum_n \delta_{nk}\,\varphi_n(\alpha) = \varphi_k(\alpha)
\end{align}

and

\begin{align}
\{\hat c_k,\hat\varPsi^\dagger(\alpha)\}
&= \sum_n\sum_m\left(\mathbf S^{-1}\right)_{nm}\varphi_m^*(\alpha)\,\{\hat c_k,\hat c_n^\dagger\}\\
&= \sum_m \delta_{km}\,\varphi_m^*(\alpha) = \varphi_k^*(\alpha)
\end{align}

We find the three anti-commutators:

\begin{empheq}[box=\fbox]{align}
\{\hat\varPsi(\alpha),\hat\varPsi^\dagger(\alpha')\} &= \delta(\alpha-\alpha')\\
\{\hat\varPsi(\alpha),\hat c_k^\dagger\} &= \varphi_k(\alpha)\\
\{\hat c_k,\hat\varPsi^\dagger(\alpha)\} &= \varphi_k^*(\alpha)
\end{empheq}

\subsection*{The Slater Determinant in second quantization}

This (anti-)commutation relation leads to a nice definition of the
Slater determinant:

\begin{equation}\label{eq:slater-determinant-first-second-connection}
\boxed{\Phi_{\alpha_1\alpha_2\dots\alpha_N}(x_1,x_2,\dots,x_N) = \frac{1}{\sqrt{N!\det(\mathbf{S_{occ}})}}\langle\mathrm{vac}|\hat\varPsi(x_N)\cdots\hat\varPsi(x_1)\,\hat c_{\alpha_N}^\dagger\cdots\hat c_{\alpha_1}^\dagger|\mathrm{vac}\rangle}
\end{equation}

The normalization factor is the same for the Slater determinant in second and first quantization and will therefore be neglected
in the following proof, that the formula \eqref{eq:slater-determinant-first-second-connection} correctly represents the Slater determinant
in first quantization.
The first nontrivial case we find for $N=2$:

\begin{align}
\hat{\Phi}_{12}(x_1,x_2)
&= \langle\mathrm{vac}|\hat\varPsi(x_1)\hat\varPsi(x_2)\,\hat c_2^\dagger\hat c_1^\dagger|\mathrm{vac}\rangle\\
&= \Big[\langle\mathrm{vac}|\hat\varPsi(x_1)\,\varphi_2(x_2)\,\hat c_1^\dagger|\mathrm{vac}\rangle
- \langle\mathrm{vac}|\hat\varPsi(x_1)\,\hat c_2^\dagger\,\hat\varPsi(x_2)\,\hat c_1^\dagger|\mathrm{vac}\rangle\Big]\\
&= \Big[\langle\mathrm{vac}|\varphi_1(x_1)\varphi_2(x_2)|\mathrm{vac}\rangle - \langle\mathrm{vac}|\varphi_2(x_2)\,\hat c_1^\dagger\,\hat\varPsi(x_1)|\mathrm{vac}\rangle\notag\\
&\qquad - \langle\mathrm{vac}|\hat\varPsi(x_1)\,\hat c_2^\dagger\,\varphi_1(x_2)|\mathrm{vac}\rangle + \langle\mathrm{vac}|\hat\varPsi(x_1)\,\hat c_2^\dagger\hat c_1^\dagger\,\hat\varPsi(x_2)|\mathrm{vac}\rangle\Big]\\
&= \Big[\varphi_1(x_1)\varphi_2(x_2) - \langle\mathrm{vac}|\varphi_2(x_1)\varphi_1(x_2)|\mathrm{vac}\rangle\notag\\
&\qquad + \langle\mathrm{vac}|\varphi_1(x_2)\,\hat c_2^\dagger\,\hat\varPsi(x_1)|\mathrm{vac}\rangle\Big]\\
&= \Big(\varphi_1(x_1)\varphi_2(x_2) - \varphi_1(x_2)\varphi_2(x_1)\Big)
\end{align}

where the tilde $\tilde{\Phi}$ denotes that the normalization factor is neglected.

\paragraph{Induction step $N\to N+1$:}
\begin{align}
\langle\mathrm{vac}|\hat\varPsi(x_1)\cdots\hat\varPsi(x_{N+1})\,\hat c_{\alpha_{N+1}}^\dagger\cdots\hat c_{\alpha_1}^\dagger|\mathrm{vac}\rangle
&= \varphi_{\alpha_{N+1}}(x_{N+1})\,\langle\mathrm{vac}|\hat\varPsi(x_1)\cdots\hat\varPsi(x_N)\,\hat c_{\alpha_N}^\dagger\cdots\hat c_{\alpha_1}^\dagger|\mathrm{vac}\rangle\notag\\
&\quad - \langle\mathrm{vac}|\hat\varPsi(x_1)\cdots\hat\varPsi(x_N)\,\hat c_{\alpha_{N+1}}^\dagger\,\hat\varPsi(x_{N+1})\,\hat c_{\alpha_N}^\dagger\cdots\hat c_{\alpha_1}^\dagger|\mathrm{vac}\rangle\notag\\
&\quad\ \vdots\notag\\
&= \sum_{i=1}^{N+1}\varphi_i(x_{N+1})\cdot \mathrm{SD}\big(\varphi_1,\dots,\varphi_{i-1},\varphi_{i+1},\dots,\varphi_{N+1}\big)\,(-1)^{N+1-i}
\end{align}

This is just the recursive formula for a determinant if you, for
example, imagine $\varphi_i(x_{N+1})$ for $i\in\{\alpha_1,\dots,\alpha_{N+1}\}$
to be the first row and $\varphi_i(x_N)$ for $i\in\{1,\dots,N+1\}$ the
second, and so on. Then ``developing'' the determinant after the first
row gives the formula above.

This result allows us to connect first and second quantization, or
abstract representation and position representation, in the same way we
connected it for the single-particle case, where $\varphi(x) = \langle x|\varphi\rangle$.
Similarly, we find for the many-particle case the abstract general
$N$-electron state:

\begin{equation}
\boxed{|\varphi\rangle = \sum_{\alpha_1<\dots<\alpha_N} a_{\alpha_1,\dots,\alpha_N}\,\hat c_{\alpha_N}^\dagger\cdots\hat c_{\alpha_1}^\dagger|\mathrm{vac}\rangle}
\end{equation}

and applying the bra

\begin{equation}
\boxed{\langle x_1,\dots,x_N| = \langle\mathrm{vac}|\hat\varPsi(x_1)\cdots\hat\varPsi(x_N)}
\end{equation}

to the abstract $N$-electron state gives the position
representation / first quantization: a linear combination of Slater
determinants.